\section{Measurement discrimination, geometry, and learning}
\label{sec:matrixcomparison76}
The present binary interface retains the ordered classical outcome
and an arbitrary external reference, but no postmeasurement quantum
system from the device. This fixes the comparison with the
measurement-discrimination literature. The projective theorem of
Pucha\l a and collaborators \cite[Corollary~1 and Theorem~2]{PPKKO72}
gives an exact multiple-use formula with ancillary access. It is
used in the retained observable-readout theorem. The restricted
ancillary interface in Fiur\'a\v sek--Mi\v cuda \cite{FM75} leads
to a different comparison of adaptive and parallel probes.

Krawiec--Pawela--Pucha\l a \cite[Section~4]{KPP76} construct a pair
of nine-outcome qutrit rank-one POVMs that admit perfect adaptive
discrimination with two calls and no perfect discrimination by
any finite parallel experiment. Datta--Biswas--Saha--Augusiak
\cite[Theorems~3--5]{DBSA76} study entanglement-assisted
single-use discrimination, including perfect constructions and a
qubit entanglement advantage. Their qubit construction has three
outcomes. These are direct antecedents for general POVM
discrimination with a retained reference. They do not supply a
full binary-effect family covering law. Our comparison
\eqref{eq:matrixadaptivity76} permits an adaptive advantage and
controls its size uniformly in $N$ for fixed input dimension.

The differential argument in Lemma~\ref{lem:matrixpath76} belongs
to the established horizontal-gauge and channel-extension method
\cite{FI76,DKG67,DM76,KGAD76}. Yuan--Fung give the parallel path estimate
\cite[Section~III, equation~(26), PDF version~3]{YF76};
Sieniawski--Demkowicz-Dobrza\'nski integrate adaptive bounds
\cite[Section~IV.A]{SD76}. The binary Sylvester choice here gives the
specific regularized variance $M(I-M)$, which can be integrated
explicitly and matched by compressed nonadaptive tests. The
inverse-Sylvester expressions of Bures geometry
\cite{Dittmann76} and its spectral integration
\cite{SZGeom76} are related methods. Our form acts on the effect
displacement $H$ with variance $V=M(I-M)$; it is not the ordinary
Bures pullback under $M\mapsto V$, whose differential would be
$H-MH-HM$. The operational-ball lemma and the nested endpoint
integral are both required for Theorem~\ref{thm:matrixcover76}.

Mele--Bittel \cite[Theorem~I.1]{MB67} prove the
fixed-confidence channel-learning bound
$O(d_{\rm in}d_{\rm out}k\epsilon^{-2})$ for Kraus rank at
most $k$ and $0<\epsilon<1$. Their binary specialization
\cite[Lemmas~III.7--III.8 and Corollary~III.9]{MB67} gives
\[
 \|\mathcal M_E-\mathcal M_F\|_\diamond
     =2\|E-F\|_{\rm op},\qquad
 M=1024d^2\epsilon^{-2}
       +O(d\epsilon^{-2}\log(1/\eta)),
\]
with $\|\widehat E-E\|_{\rm op}\le\epsilon$ and probability
at least $1-\eta$, provided $\eta>4e^{-4d^2}$.
The channel-rank bound is $k=2d$; no effect-rank assumption
is imposed. Their channel calls prepare independent maximally
entangled probes in parallel, followed by collective processing
of the resulting Choi states. The learner has no measurement-
environment access. For arbitrary confidence,
\cite[Remark~III.4]{MB67} gives the repetition bound
$O(d^2\epsilon^{-2}\log(1/\eta))$.
The dimension lower bound matches at fixed accuracy; it is
distinct from their separate accuracy lower bound.

Zambrano--Ramos-Calderer--Kueng \cite[Definition~1]{ZRK76}
use the worst-input one-use outcome total variation, denoted
$d_{\rm op}$. Their Theorem~2, equation~(13), gives
\[
 M\ge
 \frac{8(d^2+\epsilon(d^2+1)/6)}{\epsilon^2}
 \log\frac{2^{L+1}9^{2d}}{\eta}
\]
as a sufficient sample size for error $\epsilon$ and
confidence $1-\eta$, using known probes from a global
$2$-design and classical least-squares reconstruction followed
by projection onto the POVM body. Calls are nonadaptive and
single-copy, without retained references or entanglement
between calls. The result holds for arbitrary effect ranks
and $L$ outcomes. Their Theorem~9 gives the corresponding
$\Omega((d^3+d^2L)\epsilon^{-2})$ lower bound at constant
confidence within that access model. For $L=2$,
$d_{\rm op}(E,F)=\|E-F\|_{\rm op}$.

Thus both works already provide a common estimator of an
unknown binary measurement. To compare their stated guarantees
with the present loss, the hybrid inequality gives
\[
 d_N(\mathcal M_E,\mathcal M_{\widehat E})
      \le 2N\|E-\widehat E\|_{\rm op}.
\]
Substituting $\epsilon=\delta/(2N)$ into either cited
binary guarantee gives a sufficient training budget
$O_d(N^2\delta^{-2}\log(1/\eta))$.
This substitution certifies a sufficient budget; it makes no converse
assertion about those estimators under the future-use loss.
Theorem~\ref{thm:matrixlearning77} proves
$\Theta_d(N\delta^{-2}\log(1/\eta))$ directly for $d_N$,
uniformly over the complete binary matrix effect body for each fixed
$d$, for
$N\ge1$, $0<\delta\le\delta_0$, and $0<\eta\le1/8$.
Its training procedure chooses entangled input blocks of length
at most $N$ using the observed record, and returns one legal
classical estimate. The lower bound allows arbitrary retained
references, adaptive controls, and bounded public stopping.
The explicit upper bound is $Cd^4N\delta^{-2}\log(d/\eta)$;
its dimension dependence is not asserted to be sharp. Calibration of
both support endpoints and noise-adapted compression learning are
what permit the statistical theorem to share the full matrix scope
of the metric theorem.
Multiscale phase estimation and robust phase recovery are
established techniques \cite{HB76,KLY76}. Our proof supplies its
own small-cone argument and conditional confidence allocation;
it does not invoke the numerical unwrapping threshold corrected
by the erratum to \cite{KLY76}. Randomized signs remove the
unknown bias exactly, and empirical amplitudes choose the
entangled block length without a supplied noise parameter.
The qubit subprocedure is common to all unknown qubit effects.
The calibrated compression construction extends common learning to
the full matrix body in every fixed dimension. The pair-dependent tests used in the matrix
metric comparison and the supplied-target exact encoder retain
their separate roles.

The deterministic rational construction in
Theorem~\ref{thm:matrixcodec77} adds a finite exact selection rule
to the optimal-order covering theorem. Its exhaustive computation is
accounted for separately from the payload; no efficient high-dimensional
encoder or optimal arithmetic complexity is inferred.

All matrix entropy constants are for fixed dimension and small
error. The proof includes every rank and spectral multiplicity
of the binary effect body. Multi-outcome POVMs, residual quantum
outputs, sharp growing-dimension minimax and entropy laws, and full submaximal-error
matrix coverings require additional arguments. The inherited
full-error preparation and fixed-readout results retain their
own stated ranges.
